\documentclass[a4paper,12pt]{report}
\usepackage[utf8]{inputenc}
\usepackage{graphicx}
\usepackage{float}
\usepackage{tabularx}
\usepackage{makecell}
\usepackage{titlesec}
\usepackage{fancyhdr}
\usepackage{lastpage}
\usepackage{xurl}
\usepackage{hyperref}
\usepackage{geometry}
\usepackage{color}
\usepackage{microtype}
\usepackage{enumerate}
\usepackage{listings}
\usepackage[table,dvipsnames]{xcolor}
\setcounter{tocdepth}{5}
\setcounter{secnumdepth}{5} %SONO I DUE COMANDI PER AVERE LE SUBSUBSECTION NUMERATE E PRESENTI NELL'INDICE
\renewcommand{\contentsname}{Indice} %QUESTO SERVE PER AVERE L'INDICE CON IL NOME CHE VOGLIO IO
\titleformat{\chapter}[display]
{\normalfont\bfseries}{}{0pt}{\LARGE} %QUESTO SERVE PER AVERE SOLO IL NOME DEL CAPITOLO CHE VOGLIO IO
\titlespacing*{\chapter}{0cm}{0cm}{0.2cm}
\geometry{
	left=20mm,
	right=20mm,
}
\fancypagestyle{plain}{
	\fancyhf{}
	\lhead{\includegraphics[width=3cm]{../immagini/minilogo.jpg}}
	\chead{}
	\rhead{\fontsize{12}{10}Piano di progetto}
	\lfoot{}
	\cfoot{\thepage\ di \pageref*{LastPage}}
	\rfoot{}
}

\definecolor{atomlightorange}{rgb}{0.88,0.76,0.55}
\definecolor{atomdarkgrey}{RGB}{59,62,75}

\usepackage{tikz}

% set listings
\lstset{%
    basicstyle=\footnotesize\ttfamily\color{atomlightorange},
    framesep=20pt,
		belowskip=10pt,
		aboveskip=10pt
}

% add frame environment
\usepackage[%
    framemethod=tikz,
    skipbelow=8pt,
    skipabove=13pt
]{mdframed}
\mdfsetup{%
    leftmargin=0pt,
    rightmargin=0pt,
    backgroundcolor=atomdarkgrey,
    middlelinecolor=atomdarkgrey,
    roundcorner=6
}

\usepackage{etoolbox}% >= v2.1 2011-01-03
\BeforeBeginEnvironment{lstlisting}{\begin{mdframed}\vspace{-0.7em}}
\AfterEndEnvironment{lstlisting}{\vspace{-0.5em}\end{mdframed}}

% needed for \lstcapt
\def\ifempty#1{\def\temparg{#1}\ifx\temparg\empty}

% make new caption command for listings
\usepackage{caption}
\newcommand{\lstcapt}[2][]{%
    \ifempty{#1}%
        \captionof{lstlisting}{#2}%
    \else%
        \captionof{lstlisting}[#1]{#2}%
    \fi%
    \vspace{0.75\baselineskip}%
}

\hypersetup{
    colorlinks=true,
    linkcolor=black,
    filecolor=black,
    urlcolor=blue,
		citecolor=black,
}

\pagestyle{plain}

\makeatletter
	 \def\thebibliography#1{\chapter*{Bibliografia\@mkboth
		 {Bibliografia}{Bibliografia}}\list
		 {[\arabic{enumi}]}{\settowidth\labelwidth{[#1]}\leftmargin\labelwidth
 \advance\leftmargin\labelsep
 \usecounter{enumi}}
 \def\newblock{\hskip .11em plus .33em minus .07em}
 \sloppy\clubpenalty4000\widowpenalty4000
 \sfcode`\.=1000\relax}
	 \makeatother


\begin{document}

\makeatletter
\begin{titlepage}
    \begin{center}
    \vspace*{-5,0cm}
    \author{Jawa Druids}
    \title{Piano di progetto}
    \date{} %LASCIARE QUESTO CAMPO VUOTO, SE LO TOLGO STAMPA LA DATA CORRENTE
    \includegraphics[width=0.7\linewidth]{../immagini/DRUIDSLOGO.jpg}\\[4ex]
    {\huge \bfseries  \@title }\\[2ex]
    {\LARGE  \@author}\\[50ex]
    \vspace*{-8,0cm}
    \begin{table}[H]
        \centering
        \begin{tabular}{c|c}
            \textbf{Versione} & x.x.x \\
            \textbf{Data approvazione} & xx-xx-xxxx\\
            \textbf{Responsabile} & Nome Cognome\\
            \textbf{Redattori} & Andrea Dorigo \\
            \textbf{Verificatori} & Nome Cognome \\
        %MAKECELL SERVE PER POI ANDARE A CAPO ALL'INTERNO DELLA CELLA
            \textbf{Stato} & Redazione in corso\\
            \textbf{Lista distribuzione} & \makecell{Jawa Druids \\ prof. Tullio Vardanega \\ prof. Riccardo Cardin \\ Sync Lab}\\
            \textbf{Uso} & Esterno
        \end{tabular}
    \end{table}
					\vspace{.4cm}
	\hfill \break
    \fontsize{17}{10}\textbf{Sommario}\\
		\vspace{.3cm}
    Il presente documento contiene la pianificazione delle attività del gruppo Jawa Druids atte al soddisfacimento del capitolato \normalsize\textit{GDP: Gathering Detection Platform} di Sync Lab.
    \end{center}
\end{titlepage}
\makeatother

\def\myformat#1{
	\centering\huge#1
}
\hfill \break
\section*{\myformat{Registro delle modifiche}}
{\rowcolors{2}{Apricot!90!Bittersweet!20}{Bittersweet!70!Apricot!40}

\begin{table}[H]
	\centering
	\begin{tabular}{|c|c|c|c|c|}
		\hline
		\rowcolor{Melon} 
		\textbf{Modifica} & \textbf{Autore} & \textbf{Ruolo} & \textbf{Data} & \textbf{Versione}\\
		\hline
		\textit{Prima stesura bozza} & Andrea Dorigo & \textit{Responsabile} & 26-11-2020 & v0.0.1 \\
		\hline
	\end{tabular}
\end{table}
\tableofcontents{}
\chapter{Introduzione}

Lo scopo di questo documento è la stesura di un elenco di linee guida ed esempi che i componenti sono incoraggiati a seguire per migliorare l'efficacia della collaborazione.
A differenza delle Norme di Progetto, questo documento \`{e} redatto in un linguaggio pi\`{u} informale per facilitarne la comprensione, con spezzati di codice a scopo esemplificativo e riferimenti esterni sulle best practices da seguire.\\
Il documento pu\`{o} essere soggetto a modifiche ed aggiunte per tutta la durata del progetto.

\begin{thebibliography}{}
    % \bibitem{GitFlowCheatseet}
    % Daniel Kummer.
    % \textit{Git flow cheatsheet - efficient branching using git-flow by Vincent Driessen}.\\
		% \url{https://danielkummer.github.io/git-flow-cheatsheet/}.
		% \bibitem{GitFlowWorkflowTutorial}
		% Atlassian Bitbucket.
		% \textit{Git flow workflow tutorial}.\\
		% \url{https://www.atlassian.com/git/tutorials/comparing-workflows/gitflow-workflow}.
		% \bibitem{GitCommitGuidelines}
		% Cameron McKenzie, TechTarget.
		% \textit{How to write a Git commit message properly with examples}.\\
		% \url{https://www.theserverside.com/video/Follow-these-git-commit-message-guidelines}.
		% \bibitem{TabellaDelleOreLavorative}
		% Jawa Druids, 2020.
		% \textit{Tabella delle Ore Lavorative}.\\
		% \url{https://docs.google.com/spreadsheets/d/12esX1ISWQOKM-fjuHTLmAzRN0cWltksn7eGiPsBHBI0/edit?usp=sharing}.
		% \bibitem{TrelloJawaDruids}
		% Jawa Druids, 2020.
		% \textit{Trello - dashboard Jawa Druids}.\\
		% \url{https://trello.com/b/hIEOGbE9/jawadruids}.
\end{thebibliography}


\end{document}
